\section{Points, comparison maps, and central extensions}
\label{sec:D}
\providecommand{\C}{\mathbb C}
\providecommand{\PSL}{\operatorname{PSL}}

We use the final signed modular scheme of Part B, not an earlier
twisted-conjugacy parameter space.  Thus
\[
 \rho=\begin{pmatrix}0&1\\-1&1\end{pmatrix},\qquad
 \sigma=\begin{pmatrix}0&1\\-1&0\end{pmatrix},
 \quad A_R=R[r]/(r^2-r+1),\quad B_R=R[i]/(i^2+1).
\]
The norm forms are \(N_A(d+cr)=c^2+cd+d^2\) and
\(N_B(u+ti)=u^2+t^2\).  The signs \(e,e'\) below are respectively
\(\varepsilon _2,\varepsilon _3\), and \(\chi=\det U\).
This dictionary is the one in \S48, folios 542--543,
scans 479--481, especially (3) and (12bis).  It matters that
\(\mu _2(\widehat{\Z})=\prod_p\{\pm1\}\) whereas the constant group
\(C_2=\{\pm1\}\) has only two elements.  We distinguish these groups
throughout.

\subsection{Integral, local, and adelic points}

\begin{statement}{D.1}{Proved}{An explicit group law for the requested points}
Let \(R=\Z,\Z_p\), or \(\Q_p\).  Define
\[
 H_R=\{(a,z,e,e'):a\in A_R^\times,\ z\in B_R^\times,\
 e,e'\in\{\pm1\},\ N_A(a)=ee'N_B(z)\}.
\]
Put \(\kappa_e(z)=z\) for \(e=1\) and \(\kappa_e(z)=\bar z\)
for \(e=-1\).  Then the group operation is
\[
 (a,z,e,e')(a_1,z_1,e_1,e'_1)
 =(aa_1,z\kappa_e(z_1),ee_1,e'e'_1).
 \tag{D.1}
\]
There is an isomorphism of groups, and of topological groups in the
local cases,
\[
 M(R)\simeq \SL_2(R)\rtimes H_R, \qquad
 h(S)=D_{\chi(h)}S D_{\chi(h)}^{-1},\quad
 D_\mu=\operatorname{diag}(\mu,1),\quad
 \chi(h)=eN_B(z).
 \tag{D.2}
\]
For \(\Z\) this specializes to
\[
 M(\Z)\simeq C_6\times\bigl(\SL_2(\Z)\rtimes D_8\bigr),
 \tag{D.3}
\]
where \(D_8=C_4\rtimes C_2\) has order eight and acts on
\(\SL_2(\Z)\) through its quotient \(C_2\), by conjugation with
\(\operatorname{diag}(-1,1)\).
\end{statement}

\begin{proof}
The final equations give \(a=d+c\rho\) and
\[
 b=\begin{pmatrix}u&t\\-et&eu\end{pmatrix}
   =(u+et\sigma)\operatorname{diag}(1,e).
\]
Thus \(z=u+eti\) has norm \(u^2+t^2\).
Conjugation by \(\operatorname{diag}(1,e)\) acts on \(B_R\)
as \(\kappa_e\); this proves (D.1).
The remaining equation is
\(\det U=eN_B(z)=e'N_A(a)\), exactly the equation defining \(H_R\).
For fixed \(h\), write \(U=SD_{\chi(h)}\), uniquely with
\(S\in\SL_2(R)\).  Multiplication of the matrices proves (D.2);
both maps use polynomial operations and inversion of units, so are
continuous with continuous inverses.

Over \(\Z\), the two positive definite norms are positive units and
therefore equal to one.  The solutions for \(a\) are the six powers
of \(r\); those for \(z\) are \(1,i,-1,-i\).
The signs satisfy \(e'=e\).  Conjugation inverts \(i\), so
\(H_\Z=C_6\times(C_4\rtimes C_2)\), and \(\chi=e\).
This proves the asserted integral group, including its action.
\end{proof}

\begin{statement}{D.2}{Proved}{All local determinant and sign restrictions}
Set \(A_R^1=\ker N_A\) and \(B_R^1=\ker N_B\).
For \(R=\Z_p,\Q_p\) the description in D.1 has the exact sequence
of topological groups
\[
 1\longrightarrow A_R^1\times B_R^1
 \longrightarrow H_R
 \xrightarrow{(\chi,e,e')}D_R\longrightarrow1,
 \quad
 D_R=\{(\mu,e,e'):e\mu\in N_B(B_R^\times),\
                         e'\mu\in N_A(A_R^\times)\}.
 \tag{D.4}
\]
Here \(D_R\) has the multiplication inherited from
\(R^\times\times C_2^2\).  For integral local points the norm images are
\[
\begin{array}{c|cc}
 &N_A(A_{\Z_p}^{\times})&N_B(B_{\Z_p}^{\times})\\ \hline
p=2&\Z_2^\times&1+4\Z_2\\
p=3&1+3\Z_3&\Z_3^\times\\
p\ge5&\Z_p^\times&\Z_p^\times .
\end{array}                                                   \tag{D.5}
\]
Consequently \(\mu\) is any \(p\)-adic unit; at \(2\) one must
have \(e=\chi_4(\mu)\), at \(3\) one must have
\(e'=\left(\frac{\mu\bmod3}{3}\right)\), and there are no other
sign restrictions.  Here \(\chi_4(1)=1,\chi_4(3)=-1\).
The norm-one kernels themselves have the following
noncanonical topological decompositions:
\[
 \begin{array}{c|cc}
 &A_{\Z_p}^{1}&B_{\Z_p}^{1}\\ \hline
 p=2,3&C_6\times\Z_p&C_4\times\Z_p\\
 p\ge5&C_{p-(\frac{-3}{p})}\times\Z_p&
             C_{p-(\frac{-1}{p})}\times\Z_p .
 \end{array}                                                   \tag{D.5a}
\]

For \(\Q_p\), write \(\mu=p^r v\) with \(v\in\Z_p^\times\).
The condition on \(r\) is
\[
 r\in 2\Z \quad\hbox{unless }p\equiv1\pmod {12};
 \qquad r\in\Z \quad\hbox{if }p\equiv1\pmod {12}.               \tag{D.6}
\]
The restrictions at \(2,3\) are the same as above with \(v\) in
place of \(\mu\), and at other primes the two signs are arbitrary.
Equations (D.1)--(D.6), including the explicitly specified
quadratic norm-one groups, compute \(M(\Z_p)\) and \(M(\Q_p)\)
as groups, not just their determinant images.
\end{statement}
\begin{proof}
The map in (D.4) has the displayed kernel.  Its image is exactly
the displayed set: choose independently elements with the required
two norms.  This also proves surjectivity without assuming a section.

A split quadratic algebra has surjective norm on units.  In an
unramified quadratic extension, the residue-field norm is
surjective.  Norms on successive principal-unit quotients are the
residue trace, which is surjective; lifting successively and using
completeness proves surjectivity on all units.  For \(A\) the
unramified or split cases are \(p\ne3\), and for \(B\) they are
\(p\ne2\).  At \(3\), the nonzero values of
\(c^2+cd+d^2\) reduce to \(1\).  Conversely every unit \(w\equiv1\)
modulo \(3\) is \(1+d+d^2\): Hensel's lemma applies at \(d=0\)
to \(d^2+d+1-w\), whose derivative is \(1\) modulo \(3\).
At \(2\), a unit sum of two squares is \(1\) modulo \(4\).
If \(w\equiv1\pmod8\), it is a square in \(\Z_2^\times\);
if \(w\equiv5\pmod8\), \(w-4\) is such a square and
\(w=(\sqrt{w-4})^2+2^2\).  This proves (D.5).

In an unramified quadratic field the valuation of a norm is even;
all even valuations occur.  In the two ramified fields the
uniformizers \(1+i\) and \(\sqrt{-3}\) have norms \(2\) and \(3\).
Thus their norm images have arbitrary valuations and the unit
restrictions just proved.  For \(p\ge5\), the two algebras split
respectively when \(p\equiv1\pmod3\) and \(p\equiv1\pmod4\).
Both split precisely when \(p\equiv1\pmod{12}\).
At \(2\), \(A\) is unramified; at \(3\), \(B\) is unramified.
This proves (D.6) and all the sign assertions.
For (D.5a), logarithm on a sufficiently deep principal-unit
subgroup identifies the norm-one subgroup with the kernel
of the trace, a free \(\Z_p\)-module of rank one.
The compact abelian group therefore has an open copy of
\(\Z_p\).  The structure theorem for finitely generated
abelian pro-\(p\) groups splits its free rank-one part from
its torsion.  The torsion consists of the roots of unity of
norm one.  At the split odd primes there are \(p-1\);
in an unramified quadratic field there are \(p+1\).
At \(2\), the unramified field \(\Q_2(\sqrt{-3})\) has norm-one
roots \(\mu_6\), and the Gaussian field has \(\mu_4\).
At \(3\), the ramified Eisenstein field has \(\mu_6\)
and the unramified Gaussian field has norm-one roots
\(\mu_4\).  This proves the table.
Over \(\Q_p\), a field norm-one element is a unit,
so these kernels do not change; in the split case the
kernel is instead \(\Q_p^\times=\Z\times\Z_p^\times\).
The exhaustive norm computations modulo \(4,8,16,9,27,25,49\)
in \texttt{code/D\_gt.py}, recorded in \texttt{D\_gt.out},
give independent finite checks, including both ramified primes.
\end{proof}

\begin{statement}{D.3}{Proved}{Integral and finite adelic points}
For the affine finite-type group scheme \(M/\Z\),
\[
 M(\widehat{\Z})=\prod_p M(\Z_p),\qquad
 M(\mathbb A_f)=\prod\nolimits_p' M(\Q_p),
 \tag{D.7}
\]
as topological groups; the restricted product is relative to the
subgroups \(M(\Z_p)\).  The group laws and local factors are those
of D.1--D.2.  This answers the adelic question of \S49,
folio 576, scan 546, in the final modular normalization.
\end{statement}
\begin{proof}
An \(R\)-point is a tuple of finitely many coordinates satisfying
finitely many polynomial identities, with the specified inverse of
the norm.  A point over a product of rings is therefore exactly a
product of points.  The finite adeles are the union, over finite
sets \(S\) of primes, of
\(\prod_{p\in S}\Q_p\times\prod_{p\notin S}\Z_p\).
All coordinates of a point, including its inverse norm, belong
to one such ring, because there are finitely many of them.
Conversely a restricted product of local points defines these
coordinates in \(\mathbb A_f\).  These inverse maps respect the
group law and the defining product and restricted-product
topologies.
\end{proof}

\subsection{What the congruence comparison proves, and what it does not}

We write \(G_{\Q}\), \(\widehat{\mathrm{GT}}\), and
\(\widehat{\mathrm{GT}}_0\) separately.  The last group means the
continuous automorphisms described by pairs
\((\lambda,f)\in\widehat{\Z}^{\times}\times
[\widehat F_2,\widehat F_2]\) satisfying Drinfeld's relations
(I),(II); \(\widehat{\mathrm{GT}}\) also satisfies the pentagon.
We do not identify either with \(G_\Q\).
The expressions in the manuscript initially yield maps modulo
central signs (\S48, folio 543, scan 481, (24)--(25)).
A representative in an unquotiented matrix group is extra data.

\begin{statement}{D.4}{Proved}{The typed congruence map and an obstruction to openness}
Let \(\mathscr S_K\) be a group of genuinely lifted, normalized
loop solutions attached to
\(K=G_\Q,\widehat{\mathrm{GT}}\), or
\(\widehat{\mathrm{GT}}_0\), with its chosen comparison to the
profinite modular group and the triple coordinates of scan 479.
Suppose its geometric words have a genuine continuous matrix
realization extending the congruence homomorphism
\[
 q:\widehat{\SL_2(\Z)}\longrightarrow\SL_2(\widehat{\Z})
 \tag{D.8}
\]
and its multiplier is evaluated by the specified character
to \(\widehat{\Z}^{\times}\).
Evaluating the three coordinates then defines a continuous
homomorphism \(\Phi_K:\mathscr S_K\to M(\widehat{\Z})\).
This conditional construction is not an assertion that those
lift and realization data have been constructed on each bare \(K\).
If the unlifted comparison is only
\(\bar\Phi_K:K\to M(\widehat{\Z})/\Sigma\), its canonical
replacement is the pullback group
\[
 K^\sharp=K\times_{M(\widehat{\Z})/\Sigma}M(\widehat{\Z}),
 \tag{D.9}
\]
not a homomorphism from \(K\) to \(M(\widehat{\Z})\).
A homomorphism of the latter kind exists precisely when the
pullback extension has a continuous section.

For every lift that preserves the two original global signs,
the image lies in
\[
 M(\widehat{\Z})_{\mathrm{diag}}
 =\{m:(e_p(m),e'_p(m))\hbox{ is independent of }p\}.
 \tag{D.10}
\]
This is a closed subgroup of infinite index with empty interior.
Thus the proposed openness in the full signed
\(M(\widehat{\Z})\) is false whenever that comparison is defined
with the manuscript's global signs.  The same obstruction applies
to \(\widehat{\mathrm{GT}}\) and \(\widehat{\mathrm{GT}}_0\).
It also applies after quotienting by the manuscript's finite
central scalar-sign group \(\Sigma\), since the two characters
are trivial on \(\Sigma\).
It does not decide openness in an appropriately restricted target.
\end{statement}
\begin{proof}
The congruence homomorphism is induced by the compatible reduction
maps on the dense discrete group.  All identities defining a
lifted normalized triple are word and determinant identities,
hence are preserved by reduction.  The matrix triple then satisfies
the equations of \(M\).  Products of triples reduce to products,
and reduction is continuous.  This proves the asserted map
on the specified source, without asserting that the required
lift choices exist for every unlifted source.

The fibre product (D.9) has its usual group law and projection;
a section is exactly a compatible selection of a representative
for each element of \(K\).  This proves the lifting assertion.
The signs of a genuine profinite solution are elements of the
single constant group \(C_2\).  Applying them in every congruence
coordinate embeds them diagonally.  This proves (D.10).
For every \(p\ge5\), D.2 shows that \(M(\Z_p)\to C_2^2\)
is surjective.  An open subgroup of a product contains the full
factors outside a finite set, whereas (D.10) permits no independent
change of a sign at one of those primes.  It has empty interior;
its quotients obtained by recording signs at arbitrarily many
primes have unbounded order, so its index is infinite.

For an explicit point added by the scheme, let \(s\in\Z_5\)
be the square root of \(-1\) congruent to \(2\) modulo \(5\).
Take at \(5\)
\[
       a=sI,\qquad b=I,\qquad U=I,\qquad e=1,\ e'=-1,
\]
and take the identity at every other prime.  Its norm equation
is \(1=(-1)(-1)\), so D.3 gives a point of \(M(\widehat{\Z})\).
Its second sign is not diagonal, so it comes from no
global-sign-preserving solution.  The existence of \(s\) follows
from Hensel's lemma applied to \(s^2+1\) at \(2\) modulo \(5\).
Multiplying \(a,b,U\) by scalar signs changes neither their
determinants nor the conjugation sign \(e\), so both sign
characters descend modulo \(\Sigma\).  The nonopenness and
the extra point consequently persist in that quotient.
\end{proof}

\begin{statement}{D.5}{Proved}{Executed finite tests of relations (I) and (II)}
In \(\SL_2(\Z/n)\) put
\[
 x=\begin{pmatrix}-1&2\\0&-1\end{pmatrix},\quad
 y=\begin{pmatrix}-1&0\\-2&-1\end{pmatrix},\quad z=(xy)^{-1},
 \qquad K_n=\langle x,y\rangle .
\]
We use the convention
\[
 f(y,x)f(x,y)=1,\qquad
 f(z,x)z^m f(y,z)y^m f(x,y)x^m=1,\qquad m=(\lambda-1)/2.
 \tag{D.11}
\]
These are the displayed relations (I),(II) in
\citet[p.~19]{D_Schneps2012}.
Let \(L_n\) be generated by
\((x,y,z,y)\) and \((y,x,x,z)\) in \(K_n^4\).
Allow precisely the tuples in \([L_n,L_n]\) as simultaneous
evaluations
\((f(x,y),f(y,x),f(z,x),f(y,z))\), and require that
\(x\mapsto x^\lambda,\ y\mapsto f(x,y)^{-1}y^\lambda f(x,y)\)
descend to an automorphism of \(K_n\).
The exact enumeration is
\[
\begin{array}{c|rrrr|l|r}
n&|K_n|&|L_n|&|[L_n,L_n]|&\lambda\text{ modulus}
 &\lambda\text{ residues}&\text{solutions per residue}\\ \hline
2&1&1&1&2&1&1\\
3&24&24&8&12&1,5,7,11&3\\
4&4&4&1&4&1,3&1\\
6&24&24&8&12&1,5,7,11&3\\
12&96&96&8&12&1,5,7,11&3 .
\end{array}                                                   \tag{D.12}
\]
Every enumerated solution passes the stated automorphism test.
These are finite equation solutions, not a determination of the
images of any of the three infinite groups.
\end{statement}
\begin{proof}
The simultaneous substitution homomorphism \(F_2\to K_n^4\)
has image \(L_n\), and maps \(F_2'\) onto \([L_n,L_n]\).
Thus using this derived subgroup enforces that all four
evaluations arise from the same commutator word, rather than
choosing four unrelated matrices.  In a group generated by two
elements the derived subgroup is generated by the conjugates of
their commutator.  The program first forms \(L_n\), then this
normal closure, and checks the two matrix products in (D.11)
for every resulting tuple.  Powers depend only on \(\lambda\)
modulo twice the least common multiple of the orders of \(x,y,z\).
Finally a breadth-first traversal of the Cayley graph assigns
the proposed image to each vertex.  Conflicting assignments
reject a pair; a well-defined permutation of the vertices proves
that the endomorphism is an automorphism.  These finite procedures
are exhaustive, and \texttt{code/D\_gt.py} was executed:
\texttt{D\_gt.out} records every count and every surviving
\(f(x,y)\), not merely the totals.

For example at \(n=3,\lambda=1\), the three surviving evaluations
are \(I\),
\(\left(\begin{smallmatrix}1&1\\1&2\end{smallmatrix}\right)\),
and \(\left(\begin{smallmatrix}1&2\\2&2\end{smallmatrix}\right)\).
Compatibility with all finite quotients, and the pentagon in
the case of \(\widehat{\mathrm{GT}}\), have not been imposed.
Passing those tests is not a consequence of this enumeration.
\end{proof}

\begin{statement}{D.6}{Proved}{The congruence kernel is already visible through a simple quotient}
The kernel of (D.8) is nontrivial.  More precisely,
\(\widehat{\SL_2(\Z)}\) has a quotient \(\SL_2(\F_8)\) which
cannot factor through any \(\SL_2(\Z/n)\).
This conclusion concerns elements of the modular group; it does
not construct an element of \(\widehat{\mathrm{GT}}_0\), or a
solution of all profinite loop equations, in that kernel.
\end{statement}
\begin{proof}
Use \(\F_8=\F_2[t]/(t^3+t+1)\).  The matrices
\[
 u=\begin{pmatrix}1&t\\0&1\end{pmatrix},\qquad
 v=\begin{pmatrix}0&1\\1&1\end{pmatrix}
\]
have orders \(2,3\).
The execution in \texttt{D\_gt.out} verifies that they generate
all \(504\) elements and that the conjugacy-class sizes are
\[
 1,\ 56,\ 56,\ 56,\ 56,\ 63,\ 72,\ 72,\ 72 .
\]
No proper nonempty union of the nonidentity classes, together
with the identity, has order dividing \(504\); the program checks
every such union.  Thus the group is simple.
The modular presentation
\(\PSL_2(\Z)=C_2*C_3\) gives the required surjection.

By the Chinese remainder theorem, any congruence quotient is a
product of \(\SL_2(\Z/p^a)\).  A surjection from that product
to the nonabelian simple group \(S=\SL_2(\F_8)\) must be
surjective on one factor: images of factors are commuting normal
subgroups, so at most one is nontrivial.  The kernel of reduction
modulo \(p\) in that factor is a \(p\)-group, by the filtration
\(1+p^jX\mapsto X\bmod p\).  Its normal image in \(S\) is
trivial.  The surjection would therefore factor through
\(\SL_2(\F_p)\).  This group is generated by its upper and lower
unipotent matrices, each of order \(p\).  If
\(p\notin\{2,3,7\}\), all their images in a group of order
\(504=2^3 3^2 7\) are trivial.  For \(p=2,3,7\), the source
orders are \(6,24,336\), all less than \(504\).  No such
surjection exists.

If (D.8) were injective, compactness would make it a topological
isomorphism onto \(\SL_2(\widehat{\Z})\).
Every finite continuous quotient of the latter factors through
some \(\SL_2(\Z/n)\), contradicting the exhibited quotient.
\end{proof}

\begin{statement}{D.7}{Open}{The exact remaining comparison problem}
The following parts of D(b) remain undecided here.
For each of \(G_\Q,\widehat{\mathrm{GT}}\), and
\(\widehat{\mathrm{GT}}_0\), specify the actual normalized
comparison (with its central lift, or in its correct quotient),
and determine the kernel and image.
In particular we have not proved that the
\(\widehat{\mathrm{GT}}_0\)-image strictly contains the
\(G_\Q\)-image.  Nor have we produced a nonidentity solution of
all profinite loop equations annihilated by every congruence
comparison.

There are concrete starting instances.  The first nontrivial
finite lifting question in (D.12) is whether the tuple at
\(n=3,\lambda=1,f(x,y)=
\left(\begin{smallmatrix}1&1\\1&2\end{smallmatrix}\right)\)
extends to a pair in \(\widehat{\mathrm{GT}}_0\) with
\(\lambda=1\), and then whether it extends to
\(\widehat{\mathrm{GT}}\) or comes from \(G_\Q\).
For the requested invisible solution the first specific
noncongruence test supplied by our attack is
\(\SL_2(\F_8)\): find a loop solution whose congruence
coordinates are all trivial but whose coordinate in this
quotient is not, or prove none exists.

The partial results are the obstruction and extra point of D.4,
the exhaustive five-modulus computation D.5, and the actual
noncongruence quotient D.6.  The unproved step is extending
finite equation solutions compatibly through all finite
quotients.  Nontriviality of the kernel in D.6 alone supplies
no such extension.
\end{statement}

\subsection{The finite extension and the geometric group-scheme calculation}

In this subsection \(V=C_2^2\), with additive coordinates
\((u,v)\in\F_2^2\) corresponding to
\((\varepsilon _3,\varepsilon _2)=((-1)^u,(-1)^v)\).
The finite group in \S49, folios 602--608, scans 597--610,
is denoted \(E\).  This is an ordinary finite group of order eight.
Over a base on which \(2\) is invertible we may identify its
constant kernel \(C_2\) with \(\mu _2\).

\begin{statement}{D.8}{Proved}{Classification of the eight finite extension classes}
Over an algebraically closed field \(k\) of characteristic different
from \(2\), the central extensions of the constant group scheme
\(V\) by \(\mu _2\), with both end identifications fixed, form
\[
 H^2_{\mathrm{fppf}}(BV,\mu _2)=
 H^2(V,C_2)=\F_2\{u^2,uv,v^2\}.
 \tag{D.13}
\]
They are uniquely determined by
\[
 X^2=z^a,\qquad Y^2=z^b,\qquad [X,Y]=z^c,
 \qquad a,b,c\in\F_2,\quad z^2=1,
 \tag{D.14}
\]
where \(X,Y\) lift the ordered basis of \(V\).
The underlying groups are \(C_2^3\) once,
\(C_4\times C_2\) three times, \(D_8\) three times, and
the quaternion group \(Q_8\) once.

The manuscript's square rule in scan 602, (90), determines
\[
 E=\F_2^3,\qquad
 (s,u,v)(t,x,y)=(s+t+u(x+y),u+x,v+y).
 \tag{D.15}
\]
Thus its class is \(u(u+v)\), its group is \(D_8\), and it
is nonsplit.  The canonical complements over \(u=0\) and
\(u=v\) are not normal in \(E\).  Their intersection is trivial.
\end{statement}
\begin{proof}
The kernel is central, so the squares of lifts and their
commutator belong to \(\{1,z\}\).  Changing either lift by \(z\)
does not change these three quantities.  Conversely, putting
every element in the unique normal form \(z^sX^uY^v\) gives
a multiplication with central exponent
\(s+t+a ux+b vy+c vx\); it is associative because the added
function is bilinear.  It gives an extension for every triple,
and any extension with those three quantities has that normal
form and multiplication.  Hence all eight triples occur and
are inequivalent as extensions.  Addition of cocycles adds
the triples, proving the three-dimensional cohomology.
Since the finite group schemes in question are étale and \(k\)
is algebraically closed, they are constant; the same argument
classifies their fppf central extensions.  The description by
cohomology of the classifying stack is the usual normalized
descent description: a torsor on \(V\) with an associative
multiplication is a central extension.

If \(c=0\), (D.14) is abelian; it is \(C_2^3\) when
\(a=b=0\) and \(C_4\times C_2\) otherwise.
If \(c=1\), every noncentral element has order four exactly
when \(a=b=1\), giving \(Q_8\).  In the other three cases
there is a noncentral involution and a cyclic subgroup of
order four which it inverts, giving \(D_8\).

In (D.15), the square of \((s,u,v)\) is \((u(u+v),0,0)\).
This is precisely the manuscript's condition that order divide
two when \(\varepsilon _3=1\) or
\(\varepsilon _2\varepsilon _3=1\).
The lift over \((u,v)=(1,0)\) has order four, so no section
of \(E\to V\) exists.  The two canonical subgroups are
\(\{(0,0,v)\}\) and \(\{(0,u,u)\}\).
The commutator of lifts of \((u,v),(x,y)\) is
\((uy+xv,0,0)\).  Conjugating their nonidentity elements by
a suitable lift therefore multiplies them by \(z\), which
is outside the chosen complement.  Both are nonnormal.

The complete enumeration in \texttt{code/D\_finite.py} checks
all \(2^9\) normalized two-cochains: \(16\) are cocycles,
\(2\) are coboundaries, and the eight signatures in (D.14)
all occur.  It also checks the order distribution
\(1,2,2,2,2,2,4,4\) and the two failures of normality.
\end{proof}

\begin{statement}{D.9}{Proved}{The full geometric central-extension group of \(M\)}
For the final signed scheme over \(\C\), central extensions
by \(\mu _2\) form
\[
 \operatorname{Ext}_{\mathrm{cen}}(M_\C,\mu _2)
   =H^2_{\mathrm{fppf}}(BM_\C,\mu _2)
   \simeq\F_2^5.                                             \tag{D.16}
\]
The same answer holds for \(M(0)_\C\) and \(M[0]_\C\);
the latter is the \(Z^!\) to which scan 610 pulls back
the finite extension.
Three generators are the inflations of (D.13).
After choosing the two embeddings of \(A_\C\), a fourth is the
pullback of the squaring map on \(\Gm\) along the first
eigenvalue \(a_1:M_\C\to\Gm\); a fifth is its pullback along
\(\chi=\det U\).
In particular inflation of the manuscript's class \(u(u+v)\)
is nonzero.  Its two inflated complements are nonnormal;
their intersection, the identity-sign subgroup with the chosen
trivial lift, is normal.
\end{statement}
\begin{proof}
Choose eigenvalue coordinates
\((a_1,a_2)\) on \(A_\C^\times\) and
\((z_1,z_2)\) on \(B_\C^\times\).  In the identity component
the equations are
\[
 a_1a_2=z_1z_2=\det U=\delta .
\]
Thus
\[
 M_\C^0\simeq\Gm^2\times\GL_2,\qquad
 (a_1,z_1,U)\longmapsto
 ((a_1,\delta/a_1),(z_1,\delta/z_1),U).
 \tag{D.17}
\]
This is an isomorphism of algebraic groups.
Put \(a_*=(2\rho-I)/\sqrt{-3}\).  It has square \(I\)
and determinant \(-1\).  The elements
\[
 g_{e'}=(a_*,I,I),\qquad
 g_e=(a_*,\operatorname{diag}(1,-1),
                  \operatorname{diag}(-1,1))
 \tag{D.18}
\]
are commuting involutions with the indicated respective
negative signs.  They give a section \(V\to M_\C\).
In particular inflation from \(V\) is injective.

The coordinate ring in (D.17) is a localization of a polynomial
ring, hence is factorial and has trivial Picard group.  Its
units modulo squares are generated by \(a_1,z_1,\delta\);
constant units are squares in \(\C\).
The Kummer sequence therefore gives
\[
 H^1_{\mathrm{fppf}}(M_\C^0,\mu _2)
   =\F_2\{a_1,z_1,\delta\}.                                  \tag{D.19}
\]
Every such torsor is the pullback of squaring along a character,
so has a central group-extension structure.  That structure,
with its identity fixed, is unique: the difference of two
multiplications would be a morphism from the connected scheme
\(M_\C^0\times M_\C^0\) to \(\mu _2\), hence the constant
identity.  Equivalently the same argument may be made after
trivializing the torsors fppf-locally and then descending.
Thus (D.19) also computes central extensions of \(M_\C^0\).

Conjugation by \(g_{e'}\) is trivial.  Conjugation by \(g_e\)
fixes \(a_1,\delta\), sends \(z_1\) to \(\delta/z_1\), and
acts on \(U\) by an inner automorphism.  On the exponent
vector \((\alpha,\beta,\gamma)\) modulo \(2\) it acts by
\[
       (\alpha,\beta,\gamma)\longmapsto
       (\alpha,\beta,\beta+\gamma).
 \tag{D.20}
\]
The invariant subspace has dimension two, with basis
\(a_1,\delta\).  Any extension of \(M_\C\) restricts to an
invariant extension of \(M_\C^0\); both invariant generators
extend globally by the two character constructions in the
statement.

If an extension restricts trivially to \(M_\C^0\), its splitting
there is unique, since
\(\Hom(M_\C^0,\mu _2)=0\).  Conjugation consequently preserves
that splitting.  Quotienting by its image gives an extension
of \(V\), and the original extension is its pullback.
The kernel of restriction is therefore exactly the
three-dimensional inflated subspace.  Surjectivity onto the
two-dimensional invariant subspace proves (D.16).
Over \(\C\) the base has vanishing positive-degree cohomology
with these finite coefficients, so the central-extension
description is all of \(H^2(BM_\C,\mu _2)\).

For \(M(0)\), the identity component has coordinate ring
\(\C[a_1^{\pm1},z_1^{\pm1},\delta^{\pm1}]\).
For \(M[0]\), adjoining the upper-right entry of
\(U=\left(\begin{smallmatrix}\delta&b\\0&1\end{smallmatrix}\right)\)
adds one polynomial variable.  Both rings are factorial,
with the same three unit classes.  The section (D.18)
belongs to \(M(0)\), and the action on those unit classes
is again (D.20).  Every step of the restriction and
inflation argument therefore applies to both subgroups.

The section (D.18) detects the nonsplitting and nonnormality
proved in D.8.  The intersection of the complements is the
kernel of the sign map with its unique prescribed lift, and
the cocycle (D.15) is zero whenever one argument has both
signs zero.  This intersection is therefore normal.
\texttt{D\_components.py} checks (D.18) modulo \(7\), where
\(\sqrt{-3}=2\), and enumerates all eight characters in
(D.20), finding the four invariant ones.
\end{proof}

\begin{statement}{D.10}{Proved}{A genuine obstruction to an integral root-of-unity model}
There is no central extension of finite flat group schemes
\[
 1\longrightarrow\mu _2\longrightarrow E_2
   \longrightarrow\mu _2^2\longrightarrow1
 \tag{D.21}
\]
over \(\Z_2\), or over \(\Z\), whose generic fibre is the
dihedral extension (D.15).  In fact every central extension
of the displayed kind is commutative.
This does not prevent the constant group \(D_8\) from being
defined over \(\Z\); its integral kernel and quotient are
constant groups, not the indicated root-of-unity group schemes.
\end{statement}
\begin{proof}
The commutator of a central extension of a commutative group
is an alternating bicharacter
\(\mu _2^2\times\mu _2^2\to\mu _2\).
By diagonalizable-group duality,
\[
 \underline{\Hom}(\mu _2^2,\mu _2)=(C_2^2)_{\mathrm{const}}.
 \tag{D.22}
\]
For completeness, a homomorphism to \(\mu _2\) is a group-like
element of the Hopf algebra of \(\mu _2^2\).  The coefficient
identities for being group-like say that its coefficients
are pairwise orthogonal idempotents summing to one; on each
connected piece precisely one monomial is selected.  This
is the constant character group in (D.22).
The bicharacter is consequently a morphism
\(\mu _2^2\to(C_2^2)_{\mathrm{const}}\).

The source is connected over \(\Z_2\): its finite coordinate
algebra has reduction
\(\F_2[s,t]/(s^2,t^2)\), a local ring, and hence has no
nontrivial idempotent.  It is also connected over \(\Z\).
Indeed an idempotent would remain nontrivial over \(\Z_2\)
unless it or its complement became zero there; injectivity
of the base change of the free coordinate module then forces
the same vanishing over \(\Z\).
Thus the morphism to the constant group is constant, and its
value at the identity is the identity.  The commutator is
trivial.  A dihedral generic fibre is noncommutative, giving
the required contradiction.
\end{proof}

\begin{statement}{D.11}{Open}{Arithmetic descent and the remaining integral \(H^2\)}
D.8 classifies the relevant finite geometric extensions;
D.9 computes the complete central-extension group over \(\C\);
D.10 excludes the literal root-of-unity extension at \(2\).
We have not computed
the central-extension groups of \(M_{\Z_2}\) or \(M[0]_{\Z_2}\),
nor its restriction map to the generic fibre, nor the full
arithmetic descent of all five geometric classes to \(\Q\).
One exact first remaining instance is the kernel and image of
\[
 \operatorname{Ext}_{\mathrm{cen}}(M[0]_{\Z_2},\mu _2)
 \longrightarrow
 \operatorname{Ext}_{\mathrm{cen}}(M[0]_{\Q_2},\mu _2).
 \tag{D.23}
\]
The computations of all finite cocycles and all invariant
identity-component characters are attached in
\texttt{D\_finite.out} and \texttt{D\_components.out}.
Their limitation is precise: they classify étale geometric
extensions and characters; they do not classify flat
\(\mu _2\)-torsors on the nonreduced special fibre of \(M\).
That descent problem is not resolved by declaring the
constant dihedral group to be a root-of-unity group scheme.
\end{statement}

\subsection{Two distinct Brauer obstructions}

\begin{statement}{D.12}{Refuted}{The cup product on scan 611 is zero}
The finite extension (D.15), pulled back along the cyclotomic
characters
\[
 u=\chi_{-3},\qquad v=\chi_{-1}
       \quad\hbox{in }H^1(G_\Q,\mu _2),
\]
has class
\[
 c=u\smile(u+v)=\chi_{-3}\smile\chi_3
                  =(-3,3)=0\quad\hbox{in }\operatorname{Br}(\Q).
 \tag{D.24}
\]
It has order one, and remains zero over \(\Q(i)\).
Thus the claimed nonvanishing in \S49, folio 609,
scan 611, is false for the cup product actually written there.
The two quadratic characters belong to \(H^1\), not the
displayed \(H^2\).
An explicit full dihedral lift is the restriction map
\[
 G_\Q\twoheadrightarrow
 \operatorname{Gal}\bigl(\Q(\sqrt[4]{3},i)/\Q\bigr)\simeq D_8;
 \tag{D.25}
\]
its only ramified rational primes are \(2\) and \(3\).
\end{statement}
\begin{proof}
The normalized cocycle in (D.15) is
\((u,v),(x,y)\mapsto u(x+y)\), so its pullback is exactly
the cup product in (D.24).  The Kummer classes of the two
quadratic fields named on scan 611 are \(-3\) and \(3\);
the second is the product of the classes \(-3\) and \(-1\).
Under the Kummer-to-Brauer map their cup product is the
quaternion symbol.  Explicitly, let
\[
 I=\begin{pmatrix}0&-3\\1&0\end{pmatrix},\qquad
 J=\begin{pmatrix}0&3\\1&0\end{pmatrix}.
\]
Then \(I^2=-3\), \(J^2=3\), and \(IJ=-JI\).
The matrices \(1,I,J,IJ\) have determinant \(-36\) as four
vectors in \(M_2(\Q)\), so these assignments give an algebra
isomorphism \((-3,3)\simeq M_2(\Q)\).
This is also the norm identity \(N_{\Q(\sqrt{-3})/\Q}
(\sqrt{-3})=3\).  The splitting is checked exactly by
\texttt{D\_finite.py}, and the Hilbert symbols are checked
by \texttt{D\_number\_fields.py}.

Put \(\alpha=\sqrt[4]{3}>0\).  The polynomial \(X^4-3\) is
Eisenstein at \(3\), so \([\Q(\alpha):\Q]=4\); this field is
real, and adjoining \(i\) doubles its degree.
The automorphisms \(r(\alpha)=i\alpha,r(i)=i\) and
\(s(\alpha)=\alpha,s(i)=-i\) satisfy
\(r^4=s^2=1,\ srs=r^{-1}\), and generate its group.
The fixed field of the centre is \(\Q(i,\sqrt3)\).
On that biquadratic field, \(r\) has signs
\((u,v)=(1,0)\), whereas \(s\) has signs \((1,1)\);
these are exactly the order-four and involution fibres
specified in (D.15).  This identifies (D.25) as a full lift.
The polynomial discriminant is \(-2^8 3^3\), so the splitting
field is unramified outside \(2,3\).  Both primes ramify:
the quartic is Eisenstein at \(3\), and the field
contains the ramified quadratic extension \(\Q(i)\) at \(2\).
\end{proof}

\begin{statement}{D.13}{Proved}{The nonzero symbol and the meaning of \(C_4\) on scan 631}
The different quaternion algebra
\[
 (-1,-3)=(-3,-1)
 \tag{D.26}
\]
is ramified exactly at \(3\) and the real place.  Its Brauer
class has order two and generates \(C_2\), not \(C_4\).
The \(C_4\) in \S49, folio 619, scan 631, is the inverse
image of a transposition in the finite group of order twelve,
not the order of a Brauer class.  Pulling
\(0\to C_2\to C_4\to C_2\to0\) back along \(\chi_{-3}\)
gives precisely the nonzero class (D.26).
Consequently Lead 4 conflates two different extensions.
\end{statement}
\begin{proof}
For a normalized lift of \(C_2\) to \(C_4\), the central
two-cocycle is \((u,x)\mapsto ux\).
Its pullback is \(\chi_a^2=(a,a)=(a,-1)\);
the last equality follows from \((a,-a)=0\), proved by the
norm \(N(\sqrt a)=-a\).
Taking \(a=-3\) gives (D.26).  This differs from (D.24)
by the term \(\chi_{-3}^2\).

At infinity both entries of (D.26) are negative, so the
quaternion algebra is nonsplit.  At \(3\), the odd-prime
Hilbert-symbol formula gives \((-1/3)=-1\).
At \(2\), the formula for two odd units is
\((a,b)_2=(-1)^{(a-1)(b-1)/4}\), which is \(1\) here.
At every other finite prime both entries are units and the
odd-prime formula gives \(1\).
These computations appear in \texttt{D\_number\_fields.out}.
The class is killed by two because it is a cup product with
\(\mu _2\) coefficients, and is nonzero by its real
localization; hence its order is exactly two.
The exact local norm interpretation used here is
\citet[Chapter III, Theorem 4.4(d)]{D_Milne2020}.
The scan itself identifies \(C_4\) as a finite inverse-image
group and notes that its chosen \(\sigma\) has order four.
Nothing there asserts that a Brauer class has order four.
\end{proof}

\subsection{The class-field exercise: constant and arithmetic forms}

\begin{statement}{D.14}{Refuted}{The constant lift with imaginary quadratic sign}
The group in \S49, folio 606, scan 605, and folio 618,
scan 629, is
\[
 \widetilde G_3=
 \langle r,s\mid r^6=s^4=1,\ r^3=s^2,\
                         srs^{-1}=r^{-1}\rangle
 \simeq \operatorname{Dic}_3
 \simeq S_3\times_{C_2}C_4.                                  \tag{D.27}
\]
It has order twelve, with a unique element of order two.
For any continuous homomorphism \(q:G_\Q\to S_3\) whose
sign is \(\chi_{-3}\), there is no homomorphism
\(G_\Q\to\widetilde G_3\), with \(\widetilde G_3\) a
\emph{constant} finite group, lifting \(q\).
In particular making \(q\) surjective does not remove the
obstruction in scan 634.
\end{statement}
\begin{proof}
Every word in the presentation is \(r^j\) or \(r^js\)
with \(0\le j<6\); all twelve occur in \(C_3\rtimes C_4\),
where the generator of \(C_4\) inverts \(C_3\).
The quotient by its central involution is \(S_3\).
Pairing this quotient with the abelianization
\(\operatorname{Dic}_3\to C_4\) gives the fibre-product
description in (D.27): the pair map is injective and both
groups have order twelve.
Every element over an odd permutation has order four;
the only involution is the central kernel element.

This presentation agrees with the indicated modular quotient.
The normal subgroup of \(\SL_2(\Z/4)\) generated by
\(-\ell_0,-\ell_1\) is the four-element group
\[
 \left\{I,\
 \begin{pmatrix}1&2\\2&1\end{pmatrix},\
 \begin{pmatrix}3&0\\2&3\end{pmatrix},\
 \begin{pmatrix}3&2\\0&3\end{pmatrix}\right\}.
 \tag{D.28}
\]
The quotient has the presentation (D.27).
\texttt{D\_finite.py} enumerates its twelve cosets, verifies
the order distribution
\(1^1,2^1,3^2,4^6,6^2\), and checks that all six elements
above odd permutations have order four.

Complex conjugation \(c_\infty\in G_\Q\) has order two,
but \(\chi_{-3}(c_\infty)=-1\).
A lift would send it to an element above an odd permutation,
which would have order four.  A homomorphism cannot increase
the order of this element, giving a contradiction.
This proof applies even when \(q\) is not surjective.
\end{proof}

\begin{statement}{D.15}{Proved}{The exact class-field criterion for constant full lifts}
Let \(k\) be a field of characteristic different from two and
let \(q:G_k\to S_3\) be a continuous homomorphism with
nontrivial sign \(\chi_a\).  It lifts to the constant group
\(\operatorname{Dic}_3\) if and only if
\[
 \chi_a:G_k\to C_2\text{ lifts to }C_4
 \quad\Longleftrightarrow\quad
 (a,-1)=0
 \quad\Longleftrightarrow\quad
 a\in N_{k(i)/k}(k(i)^\times).
 \tag{D.29}
\]
Here if \(i\in k\), the norm condition means the split
quadratic algebra and is automatic.
If \(q\) is surjective, every lift is surjective.
When a lift exists its choices form a torsor under
\(H^1(G_k,C_2)\).

For \(k=\Q\), write the square class of \(a\) as a nontrivial
squarefree integer \(d\).  The criterion is
\[
 d>0,\qquad
 \text{every odd prime divisor of \(d\) is }1\pmod4 .
 \tag{D.30}
\]
Equivalently \(d\) is a sum of two rational squares.
This is a full criterion for the constant finite-group exercise;
the sign \(d=-3\) in the manuscript fails it.
\end{statement}
\begin{proof}
By (D.27), a lift is exactly a pair \((q,\eta)\) with
\(\eta:G_k\to C_4\) reducing to \(\chi_a\).
The obstruction cocycle of \(C_4\to C_2\) is
\(\chi_a\smile\chi_a\), calculated in D.13 as \((a,-1)\).
A quaternion symbol \((-1,a)\) is split exactly when \(a\)
is a norm from \(k(i)\): to see this directly, write the
quaternion algebra as \(k(i)\oplus k(i)j\) with \(j^2=a\)
and \(jz=\bar z j\).  If \(a=N(b)\), replacing \(j\) by
\(b^{-1}j\) makes its square one, and its action together
with multiplication by \(k(i)\) identifies the algebra
with \(\operatorname{End}_k(k(i))\).
Conversely a nonzero zero divisor \(z+w j\), which exists in
a split algebra, has reduced norm
\(N(z)-aN(w)=0\); necessarily \(w\ne0\), giving
\(a=N(z/w)\).  This proves (D.29).
The ratio of two lifts takes values in the central kernel
and is a character, proving the torsor assertion.
If \(q\) is onto, a proper image of a lift would have order
six and split the extension of \(S_3\), impossible because
of its involutions.  Its image is therefore all twelve
elements.

Here is the class-field-theoretic reduction over \(\Q\).
By \citet[Chapter I, Theorem 4.16]{D_Milne2020},
\(G_\Q^{\mathrm{ab}}\simeq\widehat{\Z}^{\times}\) through
the cyclotomic action.  Thus the problem is precisely whether
the quadratic Dirichlet character \(\chi_d\) has a continuous
quartic character whose square is \(\chi_d\).
Such a square is trivial on every element of order two in
\(\prod_p\Z_p^\times\).
At an odd ramified prime the local quadratic character is
the Legendre character; on \(-1\) it is
\((-1)^{(p-1)/2}\).  Hence primes \(3\bmod4\) cannot occur.
Complex conjugation gives \(d>0\).

Conversely, for each \(p\equiv1\pmod4\) dividing \(d\),
choose a generator of \(\F_p^\times\) and send it to \(i\);
this gives a quartic character whose square is the Legendre
character.  If \(2\mid d\), the remaining local character
is \(\chi_8\).  On
\(\Z_2^\times=\{\pm1\}\times\overline{\langle5\rangle}\),
the character sending \(-1\) to \(1\) and \(5\) to \(i\)
factors modulo \(16\) and has square \(\chi_8\).
Multiplying these finitely many local characters gives the
required continuous quartic character.
This proves (D.30) by class field theory, rather than merely
asserting a lifting theorem.
The equivalence with a sum of squares is the norm expression
\(N_{\Q(i)/\Q}(x+iy)=x^2+y^2\) in (D.29).
\end{proof}

\begin{statement}{D.16}{Proved}{Explicit full lifts with complete local ramification}
There are two useful corrections to the exercise on scan 634.

\emph{Constant group over \(\Q\).}
Let \(L_5\) be the splitting field of \(X^3-12X-14\).
Then
\[
 K_5=L_5\,\Q(\zeta_5),\qquad
 \operatorname{Gal}(K_5/\Q)=\operatorname{Dic}_3 ,
 \tag{D.31}
\]
has degree twelve and sign field \(\Q(\sqrt5)\).
Its ramified primes, with ramification index \(e\), residue
degree \(f\), and number \(g\) of primes, are
\[
 \begin{array}{c|rrr}
 p&e&f&g\\ \hline
 2&3&4&1\\
 3&3&4&1\\
 5&4&1&3 .
 \end{array}                                                 \tag{D.32}
\]

\emph{The imaginary sign with its arithmetic determinant action.}
Retain the order-two arithmetic automorphism \(\tau\) of
\(\operatorname{Dic}_3\) which fixes \(r\) and sends
\(s\) to \(s^{-1}\).  Then
\[
 G_{\mathrm{ar}}=\operatorname{Dic}_3\rtimes\langle\tau\rangle
 \simeq\GL_2(\Z/4)/V_4
 \simeq S_3\times_{C_2}D_8 .
 \tag{D.33}
\]
where \(V_4\) is the subgroup (D.28).
Let \(L\) be the splitting field of \(X^3-2\), and set
\[
 K=L\,\Q(\sqrt[4]{3},i)
    =\Q(\sqrt[3]{2},\sqrt[4]{3},i).
 \tag{D.34}
\]
Then \(\operatorname{Gal}(K/\Q)=G_{\mathrm{ar}}\) has order
twenty-four, its determinant character is \(\chi_{-1}\),
and its \(S_3\)-sign is \(\chi_{-3}\).
Over \(\Q(i)\) it is a genuine full
\(\operatorname{Dic}_3\)-extension.
The only ramified rational primes are \(2,3\); at both of
them \(K/\Q\) has \((e,f,g)=(12,2,1)\).
For \(K/\Q(i)\), the corresponding pairs \((e,f)\) are
\((6,2)\) at the prime over \(2\) and \((12,1)\) at the
prime over \(3\).

More generally, every surjection \(G_\Q\to S_3\) with sign
\(\chi_{-3}\) has a full lift to \(G_{\mathrm{ar}}\) with
determinant \(\chi_{-1}\), by the same compositum with
\(\Q(\sqrt[4]{3},i)\).  Over \(\Q\) this gives a
\(\operatorname{Dic}_3\)-valued \emph{twisted cocycle},
not the constant-group homomorphism excluded by D.14.
\end{statement}
\begin{proof}
The cubic \(X^3-12X-14\) has no rational root, and its
discriminant is \(1620=18^2\cdot5\).
Its splitting group is therefore \(S_3\) and its quadratic
subfield is \(\Q(\sqrt5)\).
It is Eisenstein at \(2\); after replacing \(X\) by \(X-1\)
it becomes \(X^3-3X^2-9X-3\), Eisenstein at \(3\).
The corresponding local integer rings are generated by these
uniformizers.  Only \(2,3\) can divide the index, since their
squares are the only prime squares dividing \(1620\);
the two Eisenstein index tests show that neither does.
The field discriminant is \(1620\).
The cyclic quartic field \(\Q(\zeta_5)\) has the same
quadratic subfield.  Its intersection with an \(S_3\)
extension is at most quadratic, since every abelian quotient
of \(S_3\) factors through its sign.  Thus the intersection
is exactly \(\Q(\sqrt5)\).  The compositum has the
fibre-product Galois group in (D.27), proving (D.31).

At \(2,3\) the cubic fields are totally ramified of degree
three.  In each case \(\Q_p(\sqrt5)\) is the unramified
quadratic extension, so the local splitting fields have
\((e,f)=(3,2)\).  The order of \(2\), respectively \(3\),
modulo \(5\) is four, so adjoining \(\zeta_5\) gives the
unramified extension of degree four.  The composita have
\((e,f)=(3,4)\).  At \(5\) the cubic reduces to
\((X-2)^2(X-1)\).  Its simple root lifts, and its remaining
quadratic splitting field is \(\Q_5(\sqrt5)\).
It is contained in the totally ramified quartic field
\(\Q_5(\zeta_5)\), so \((e,f,g)=(4,1,3)\).
The discriminants show there is no other ramification.
The exact irreducibility, index, factorization, Hilbert-symbol,
and quartic-character checks were executed in
\texttt{D\_number\_fields.py}; its output also reports an
independent round-two computation of the field discriminant.

For the second construction, the subgroup (D.28) is normal
also in \(\GL_2(\Z/4)\).  Conjugation by
\(\operatorname{diag}(-1,1)\) fixes the class of \(\rho\)
and inverts that of \(\sigma\).
\texttt{D\_arithmetic.py} verifies normality in all \(96\)
matrices, the \(24\) cosets, and these two conjugation
identities.  Factoring \(G_{\mathrm{ar}}\) by its normal
\(C_3\) gives \(D_8\), and factoring by the appropriate
normal subgroup of order four gives \(S_3\).
The resulting injection into their fibre product is an
isomorphism by orders; both quotients have the same parity
map to \(C_2\).  This proves (D.33).
To specify that parity without ambiguity, in
\(F=\Q(\alpha,i)\), \(\alpha=\sqrt[4]{3}\), write
\(R(\alpha)=i\alpha,\ R(i)=i\), and let \(B\) be complex
conjugation.  The common character is \(\chi_{-3}\);
it is negative on both \(R\) and \(B\).
The determinant character is \(\chi_{-1}\), negative
on \(B\) and positive on \(R\).
In the fibre product we can take
\(s=(\text{a transposition},R)\) and
\(\tau=(1,RB)\).
Thus \(\tau\) has trivial common parity, inverts \(s\),
and fixes the generator \(r=(\text{a three-cycle},R^2)\).

The splitting field \(L\) has group \(S_3\) and unique
quadratic field \(\Q(\sqrt{-3})\), by the discriminant
\(-108\).  The dihedral field \(F=\Q(\sqrt[4]{3},i)\)
also contains that quadratic field.  Since
\(\gcd([L:\Q],[F:\Q])=2\), their intersection is precisely
that field.  Therefore \(K\) has the fibre-product group
in (D.33).  The same argument works for every \(S_3\)
extension with that quadratic sign field, proving the
general full-image assertion.

We verify all local degrees.  At \(2\), \(X^3-2\) is
Eisenstein and \(\zeta_3\) gives an unramified quadratic
extension; thus \(L_2\) has \((e,f)=(3,2)\).
The polynomial \((X+1)^4-3\) is Eisenstein at \(2\);
therefore \(\Q_2(\sqrt[4]{3})\) is totally ramified of
degree four.  It cannot contain \(i\): otherwise it would
contain \(\sqrt{-3}=i\sqrt3\), which generates the
unramified quadratic extension of \(\Q_2\).
Thus \(F_2\) has degree eight, residue degree two and
ramification index four.  The local intersection is the
common quadratic field, so \([K_2:\Q_2]=24\).
The ramification indices \(3,4\) force \(e\ge12\), and
the residue degree is at least two; hence \(e=12,f=2\).

At \(3\), \((X-1)^3-2\) is Eisenstein.
Adjoining \(\zeta_3\) adds a ramified quadratic extension,
so \(L_3\) has \(e=6,f=1\).
The field generated by \(\sqrt[4]{3}\) is totally ramified
of degree four, and \(i\) is unramified quadratic; thus
\(F_3\) has \(e=4,f=2\).
Again the intersection is the common quadratic field, so
the compositum has degree twenty-four.  The same bounds
give \(e=12,f=2\).
Both primes have just one prime above them since these
local degrees equal the global degree.
All other primes are excluded by the two polynomial
discriminants.  Dividing the ramification and residue
degrees by those of \(\Q(i)/\Q\) gives the relative table.
Every Eisenstein and quadratic residue calculation used
here is recorded in \texttt{D\_arithmetic.out}.

Finally decompose the full arithmetic representation as
\(\gamma\mapsto(h(\gamma),\chi_{-1}(\gamma))\) using the
semidirect product.  Its first coordinate satisfies
\[
 h(\gamma\gamma')=
 h(\gamma)\,\tau^{\,v(\gamma)}(h(\gamma')),
 \qquad (-1)^{v(\gamma)}=\chi_{-1}(\gamma).
 \tag{D.35}
\]
This explains exactly the arithmetic twisting.  On
\(G_{\Q(i)}\) the twist is trivial, and full image in
\(G_{\mathrm{ar}}\) implies full image
\(\operatorname{Dic}_3\) on this subgroup.
\end{proof}

\subsection{The generation remark}

\begin{statement}{D.17}{Proved}{Exact modular indices and the failure of a general finite correction}
Take
\(\tau=\operatorname{diag}(-1,1)\).
For every \(n\ge1\),
\[
 \langle\rho,\sigma\rangle=\SL_2(\Z/n),\qquad
 \langle\rho,\sigma,\tau\rangle
     =\det^{-1}(\{1,-1\})\subset\GL_2(\Z/n).
 \tag{D.36}
\]
Consequently the index in \(\GL_2(\Z/n)\) is one for
\(n=1,2\), and \(\varphi(n)/2\) for \(n>2\).
The quotient is explicitly
\((\Z/n)^\times/\{\pm1\}\).
The program enumerated the generated matrix subgroups for
every \(n\le48\), obtaining the following indices:
\[
\begin{array}{r|rrrrrrrrrrrr}
n&1&2&3&4&5&6&7&8&9&10&11&12\\ \hline
\mathrm{index}&1&1&1&1&2&1&3&2&3&2&5&2\\[2pt]
n&13&14&15&16&17&18&19&20&21&22&23&24\\ \hline
\mathrm{index}&6&3&4&4&8&3&9&4&6&5&11&4\\[2pt]
n&25&26&27&28&29&30&31&32&33&34&35&36\\ \hline
\mathrm{index}&10&6&9&6&14&4&15&8&10&8&12&6\\[2pt]
n&37&38&39&40&41&42&43&44&45&46&47&48\\ \hline
\mathrm{index}&18&9&12&8&20&6&21&10&12&11&23&8 .
\end{array}                                                   \tag{D.37}
\]
Topologically,
\[
 \overline{\langle\rho,\sigma\rangle}
    =\SL_2(\widehat{\Z}),\qquad
 \overline{\langle\rho,\sigma,\tau\rangle}
    =\det^{-1}(\{\pm1\}_{\mathrm{diag}})
       \subset\GL_2(\widehat{\Z}).
 \tag{D.38}
\]
The second subgroup has infinite index and is not open.
Thus scan 652 (\S49, folio 628 verso, lower pagination 651)
has no general
finite-group correction: in the modular special-linear
case the correction is trivial, whereas in the full
general-linear case it is infinite.
For \(n>2\), \(\tau\) itself is not in the determinant-one
group, so a phrase \(\langle\rho,\sigma,\tau\rangle
\subset G^+\) requires a definition of \(G^+\) containing
orientation reversal.
\end{statement}
\begin{proof}
The matrix \(\sigma^{-1}\rho=
\left(\begin{smallmatrix}1&-1\\0&1\end{smallmatrix}\right)\)
is an elementary unipotent; its powers give every upper
elementary matrix over \(\Z/n\).
Conjugating by \(\sigma\) gives every lower elementary
matrix.  These generate \(\SL_2\) over a local ring:
in the first column of an invertible matrix one entry is
a unit, interchange the rows if necessary, eliminate the
lower entry, and reduce to a diagonal matrix
\(\operatorname{diag}(a,a^{-1})\).
The latter is elementary because
\[
 \begin{pmatrix}1&a\\0&1\end{pmatrix}
 \begin{pmatrix}1&0\\-a^{-1}&1\end{pmatrix}
 \begin{pmatrix}1&a\\0&1\end{pmatrix}
 =\begin{pmatrix}0&a\\-a^{-1}&0\end{pmatrix};
\]
multiply by the inverse of the same expression at \(a=1\).
For \(\Z/n\), apply this argument to each prime-power
factor and use the Chinese remainder theorem to choose
elementary parameters supported on one factor.
This proves the first equality in (D.36).

The determinant of every word in the three generators is
\(\pm1\); conversely multiplying a matrix of determinant
\(-1\) by \(\tau\) gives determinant one.
The determinant map from \(\GL_2(\Z/n)\) onto
\((\Z/n)^\times\) is split by \(\operatorname{diag}(a,1)\).
This proves the second equality and index formula.
The breadth-first enumeration in \texttt{D\_finite.py}
checked both generated group orders separately, using
\[
 |\SL_2(\Z/n)|=n^3\prod_{p\mid n}(1-p^{-2})
\]
as the expected order; the output gives all forty-eight
rows, not just the indices displayed here.

Finite-level surjectivity implies density in
\(\SL_2(\widehat{\Z})\), because congruence neighbourhoods
are specified by one integer \(n\).
Adjoining \(\tau\) gives exactly the two determinant fibres
in (D.38).  The quotient
\(\widehat{\Z}^\times/\{\pm1\}_{\mathrm{diag}}\) is infinite;
the indices (D.37), or primes \(p\) of arbitrarily large
size, already show this.  A closed finite-index subgroup
of a profinite group is open, and this subgroup cannot be
open since its finite-level indices are unbounded.
The same obstruction persists in more general reductive
groups with an additional torus on which the chosen
generators act trivially.
\end{proof}

\subsection{A precise motivic construction and its necessary extra data}

For the question in \S48, folio 547, scan 489, we use
\emph{elliptic motivic representation} to mean an object of
the realization category \(\mathrm{MEM}_*\) of
\citet[Definition 6.1]{D_HainMatsumoto2020}.
Its data are a mixed Tate fibre at the specified cusp,
its weight filtration \(M\), a second filtration \(W\),
Betti monodromy, a filtered rational de Rham bundle with
integrable logarithmic connection, and comparison maps.
The residue is nilpotent, the connection preserves \(W\)
and satisfies Griffiths transversality, the \(\ell\)-adic
monodromy is \(G_\Q\)-equivariant, and the \(W\)-graded
pieces are sums of \(\operatorname{Sym}^m H(r)\).
Thus arithmetic compatibility is part of the data.
An integral version additionally specifies compatible
\(\Z_\ell\)-lattices; a rational realization alone does not
specify finite torsion.

\begin{statement}{D.18}{Known}{Specialization for elliptic torsion and unipotent fundamental groups}
Let \(E/\Q\) be an elliptic curve, with an actual rational
base point in \(E-\{0\}\) or a nonzero rational tangent
vector at \(0\) whenever a fundamental group is used.
The universal elliptic realization specializes to
continuous representations
\[
 G_\Q\longrightarrow\Aut(E[n](\overline\Q))
                     \simeq\GL_2(\Z/n),\qquad
 G_\Q\longrightarrow\GL(T_\ell E)\simeq\GL_2(\Z_\ell),
 \tag{D.39}
\]
with determinant the cyclotomic character.
Changing a basis conjugates these homomorphisms.
The universal unipotent Lie algebra over
\(\mathcal M_{1,2}\), or its tangential restriction, is a
pro-object of this realization category and specializes to
the genuine continuous action on
\(\operatorname{Lie}\pi_1^{\mathrm{un}}(E-\{0\},x)\).
The relevant results are
\citet[Example 6.6, Theorem 6.9, Corollary 6.10]{D_HainMatsumoto2020}.

For \(X=\mathbb P^1_\Q-\{0,1,\infty\}\) with tangential
base point \(\overrightarrow{01}\), its geometric
\(\Q_\ell\)-unipotent completion agrees with the rational
unipotent completion of \(F_2\) tensored with \(\Q_\ell\).
Writing \(\mathfrak p_\ell\) for its Lie algebra and
\(L^\bullet\) for the lower central series,
\[
 W_{-2m}\mathfrak p_\ell=L^m\mathfrak p_\ell,\qquad
 \operatorname{Gr}^{W}_{2m+1}\mathfrak p_\ell=0 .
 \tag{D.40}
\]
The Zariski closure of the Galois image is a negatively weighted extension of
\(\Gm\) by a prounipotent group; the weighted completion
maps onto it.  On positive graded Lie algebras its outer
image is the rationalization of the corresponding
Ihara filtration.
These assertions are the reductions of
\citet[Theorem A.6, Proposition 8.3, Theorem 8.4]{D_HainMatsumoto2001}
given below, not an identification of a profinite group
with its unipotent completion.
\end{statement}
\begin{proof}
Let \(\pi:\mathcal E\to\mathcal M_{1,1}\) be the universal
elliptic curve.  Over \(\Z[1/n]\), multiplication by \(n\)
has finite étale kernel of rank \(n^2\).
Its geometric fibre is \((\Z/n)^2\), so the finite étale
sheaf \(\mathcal E[n]\) is a representation of the
arithmetic étale fundamental group of the stack.
Pulling this sheaf back along the object
\(\Spec\Q\to\mathcal M_{1,1}\) defined by \(E\) gives
exactly \(E[n]\).  On geometric torsion points this is
the usual action obtained by applying Galois automorphisms
to their coordinates, proving (D.39).
The transition maps in \(n=\ell^r\) commute with this
pullback.  Taking their inverse limit gives \(T_\ell E\)
and proves continuity.

The Weil pairing is functorial and identifies
\(\bigwedge^2 E[n]\) with \(\mu_n\).
Applying \(\gamma\in G_\Q\) to a symplectic basis multiplies
its pairing by \(\det(\gamma|E[n])\); applying it to the
root of unity multiplies the exponent by the cyclotomic
character.  These exponents are equal.  Passing to the
inverse limit proves the determinant statement for
\(T_\ell E\).
At primes dividing \(n\), \(\mathcal E[n]\) is finite flat,
and the preceding \emph{étale} assertion is deliberately
restricted to \(\Z[1/n]\); no equality of its special-fibre
geometric point counts with \(n^2\) is being asserted.

In the cited elliptic results, \(H=R^1\pi_*\Q\) is the
basic object, and the punctured-elliptic unipotent Lie
algebra is the stated pro-object.
For each finite lower-central quotient, pull back its
compatible realizations along the chosen rational or
tangential object.  Functoriality of pullback preserves
the bracket, filtrations, comparisons, and transition
maps.  The resulting finite-dimensional Galois actions
are compatible, so their inverse limit is the action
on the pro-Lie algebra.  The rational base point gives
an actual fibre functor and hence a genuine action.
Without it the construction gives an outer action;
an identification with a different fixed fibre requires
a path and introduces precisely conjugation by that
path's geometric monodromy.
This explains the exact reduction from the universal
theorem to the representations at issue.

For the genus-zero assertion, the units \(x\) and \(1-x\)
give a \(G_\Q\)-invariant Kummer basis of
\[
 H^1_{\mathrm{\acute et}}(X_{\overline\Q},\Q_\ell(1))
       =\Q_\ell^2 .
 \tag{D.41}
\]
Indeed the geometric first cohomology has rank two, and
the two classes are independent by their residues at
\(0\) and \(1\).
The curve and the specified tangent vector have good
integral models away from \(\ell\), so the action factors
through \(G_{\Q,\{\ell\}}\), the maximal extension
unramified outside \(\ell\) and the real place.
Equation (D.41) checks the precise hypothesis of
Hain--Matsumoto Proposition 8.3, giving (D.40), the
negative weights, and the surjection from weighted
completion.  Its Theorem 8.4 then identifies the
associated positive graded outer image after tensoring
the Ihara graded groups with \(\Q_\ell\); the integral
comparison has only finite kernel and cokernel at each
degree as in its Proposition 8.2.
Finally \(\pi_1^{\mathrm{top}}(X(\C))=F_2\) is finitely
generated, which is exactly the hypothesis of Theorem A.6.
It identifies the unipotent completion of its profinite
or pro-\(\ell\) completion with the rational unipotent
completion after scalar extension.

As a finite check of the resulting free Lie algebra,
\texttt{D\_motives.py} computes the first six lower-central
graded dimensions from the Witt formula:
\(2,1,2,3,6,9\).  These are dimensions of the Lie algebra
of the geometric fundamental group, not dimensions of
the unknown full Galois image.
\end{proof}

\begin{statement}{D.19}{Refuted}{An arbitrary integral discrete representation is insufficient}
There is no construction taking every integral representation
of \(\SL_2(\Z)\) to an object of \(\mathrm{MEM}_1\) with
that Betti monodromy.  In addition, a geometric representation
alone does not determine a unique arithmetic representation
over \(\Q\).
\end{statement}
\begin{proof}
Use the discrete presentation
\(\SL_2(\Z)=\langle\rho,\sigma\mid\rho^3=\sigma^2,\
\rho^6=\sigma^4=1\rangle\).
The assignments
\[
 \rho\longmapsto B=\begin{pmatrix}0&-1\\1&-1\end{pmatrix},
 \qquad
 \sigma\longmapsto A=\begin{pmatrix}-1&0\\0&1\end{pmatrix}
 \tag{D.42}
\]
give a homomorphism to \(\GL_2(\Z)\), since
\(B^3=A^2=I\).
The cusp generator \(\sigma^{-1}\rho\) has image
\[
 AB=\begin{pmatrix}0&1\\1&-1\end{pmatrix},
 \qquad \det(TI-AB)=T^2+T-1 .
 \tag{D.43}
\]
One eigenvalue has absolute value greater than one, so
no positive power is unipotent.  Nilpotent logarithmic
residue in the definition of \(\mathrm{MEM}_1\) forces
unipotent cusp monodromy.  Thus (D.42) cannot underlie
such an object, even after a finite cover of the cusp.
All matrix identities and the characteristic polynomial
were executed in \texttt{D\_motives.py}.

For arithmetic nonuniqueness take
\[
 E:y^2=x^3-x,\qquad E^{(2)}:y^2=x^3-4x .
 \tag{D.44}
\]
They are isomorphic over \(\overline\Q\) and their
\(\ell\)-adic representations differ by \(\chi_2\).
At \(5\), direct enumeration gives
\(|E(\F_5)|=8\) and \(|E^{(2)}(\F_5)|=4\).
Their Frobenius traces are therefore \(-2\) and \(2\),
so the arithmetic representations for \(\ell\ne5\) are
not isomorphic, although the geometric fibres are.
The same argument applies to a universal family and its
constant quadratic twist: they have isomorphic geometric
local systems but different arithmetic descent.
\texttt{D\_motives.out} records this count and five more
primes, checking the twist relation at each one.
We use this example over \(\Q\); it does not assert that
both twists belong to the everywhere integral category
\(\mathrm{MEM}_1\).
\end{proof}

\begin{statement}{D.20}{Proved}{The extra data required for general arithmetic descent}
Fix a rational or tangential section of an arithmetic
fundamental-group sequence, thereby writing its group as
\(\Pi=\Gamma\rtimes G_\Q\), with continuous action of
\(G_\Q\) on the geometric profinite group \(\Gamma\).
A continuous geometric representation
\(\rho:\Gamma\to\GL(V_\ell)\) extends to a continuous
representation of \(\Pi\) if and only if one gives a
continuous homomorphism \(a:G_\Q\to\GL(V_\ell)\) satisfying
\[
 a(\gamma)\rho(h)a(\gamma)^{-1}=\rho(\gamma h)
       \quad(\gamma\in G_\Q,\ h\in\Gamma).
 \tag{D.45}
\]
Given one such \(a\), the extensions up to conjugation
by the centralizer \(C=Z_{\GL(V_\ell)}(\rho(\Gamma))\)
are classified by continuous nonabelian
\(H^1(G_\Q,C)\), where \(G_\Q\) acts on \(C\) by
conjugation through \(a\).

A construction with the motivic meaning fixed above must
add to this descent datum the mixed Tate fibre,
de Rham and Hodge structures, boundary conditions,
comparison maps, and, for torsion, stable integral
lattices.  D.18 supplies these in the two requested
geometric cases.  Neither the matrix equations of \(M\)
nor a discrete integral monodromy representation alone
provide them.
\end{statement}
\begin{proof}
Restriction of an extension \(\widetilde\rho:\Pi\to\GL(V_\ell)\)
to the specified section gives \(a\), and the conjugation
relation in a semidirect product gives (D.45).
Conversely
\(\widetilde\rho(h,\gamma)=\rho(h)a(\gamma)\)
is continuous and multiplicative precisely because
of (D.45).
For another extension \(a'\), put
\(b(\gamma)=a'(\gamma)a(\gamma)^{-1}\).
Equation (D.45) for both extensions says \(b(\gamma)\in C\);
multiplicativity says
\[
 b(\gamma\gamma')=
 b(\gamma)\,a(\gamma)b(\gamma')a(\gamma)^{-1}.
 \tag{D.46}
\]
This is exactly the continuous cocycle equation.
Conjugation by an element of \(C\) changes \(b\) by the
usual coboundary, proving the classification.
The additional realization requirements are those of
the explicit definition above, not consequences of
(D.45).  D.19 proves that omitting them changes the
assertion to a false one.
\end{proof}
